\documentclass[11pt, a4paper]{article}

\usepackage[T1]{fontenc}
\usepackage[utf8]{inputenc}
\usepackage{tgpagella}
\usepackage[spanish, es-noshorthands]{babel}
\usepackage{amsmath, amssymb, amsfonts}
\usepackage{graphicx}
\usepackage{booktabs}
\usepackage[most]{tcolorbox}
\usepackage[american]{circuitikz}
\usepackage[margin=2.5cm]{geometry}
\usepackage{fancyhdr}

\pagestyle{fancy}
\fancyhf{}
\setlength{\headheight}{16pt}
\lhead{\textbf{UCB Sede Tarija} | Ing. Mecatrónica}
\chead{Circuitos Electrónicos II (IMT-221)}
\rhead{EC1 Paralelo Diferido}
\lfoot{Elemento de Competencia 1}
\rfoot{Página \thepage}
\renewcommand{\headrulewidth}{0.4pt}

\definecolor{ucbblue}{RGB}{0, 51, 102}

\begin{document}

\begin{center}
    \Large\textbf{\textcolor{ucbblue}{UNIVERSIDAD CATÓLICA BOLIVIANA "SAN PABLO" — SEDE TARIJA}}\\
    \vspace{0.2cm}
    \normalsize Asignatura: Circuitos Electrónicos II (IMT-221) \\
    \textbf{EC1 Diferido (Paralelo) — Elemento de Competencia 1}[cite: 5]
\end{center}

\vspace{0.2cm}
\begin{tcolorbox}[colback=ucbblue!5!white, colframe=ucbblue, title=\textbf{Condiciones de la Evaluación}]
    \begin{itemize}
        \item \textbf{Tiempo:} 55 minutos[cite: 5].
        \item \textbf{Puntaje:} 50 puntos[cite: 5].
        \item \textbf{Material permitido:} calculadora científica y formulario de una hoja[cite: 5].
        \item Muestre todo el desarrollo algebraico. Defina claramente las impedancias y referencias fasoriales empleadas[cite: 5].
        \item Una respuesta aislada sin procedimiento recibe crédito limitado conforme a la clave[cite: 5].
        \item Se utilizarán fasores de \textbf{valor pico}[cite: 5].
    \end{itemize}
\end{tcolorbox}

\vspace{0.4cm}
\section*{Problema 1 — Red pasiva RLC cargada (25 puntos)[cite: 5]}
Una fuente sinusoidal ideal $\hat V_{in} = 2.0\angle0^\circ\,\text{V}_{\text{pico}}$ opera a $f = 1\,\text{kHz}$[cite: 5]. La fuente alimenta una resistencia $R_1 = 1\,\text{k}\Omega$ en serie con una red conectada entre el nodo $A$ y tierra[cite: 5]. Desde el nodo $A$ existen dos ramas en paralelo[cite: 5]:
\begin{itemize}
    \item \textbf{Rama 1:} Capacitor $C = 159.15\,\text{nF}$[cite: 5].
    \item \textbf{Rama 2:} Resistor $R_2 = 1\,\text{k}\Omega$ en serie con un inductor $L = 159.15\,\text{mH}$[cite: 5].
\end{itemize}
La salida $\hat V_{out}$ se mide \textbf{únicamente sobre el inductor}[cite: 5].

\begin{center}
    \begin{circuitikz}[scale=1, transform shape, thick]
        \draw (0,0) to[V, v=$\hat V_{in}$] (0,3) to[R, l=$R_1$] (3,3) coordinate(A) node[above]{$A$};
        \draw (A) to[C, l=$C$] (3,0) -- (0,0);
        \draw (A) -- (6,3) to[R, l=$R_2$] (6,1.5) to[L, l=$L$, v=$\hat V_{out}$] (6,0) -- (3,0);
        \draw (3,0) node[ground]{};
    \end{circuitikz}
\end{center}

\textbf{a) Modelado de impedancia (10 puntos)[cite: 5]} \\
A la frecuencia dada, determine $Z_C$, $Z_L$, la impedancia total de la rama $R_2+L$, y calcule la impedancia equivalente vista desde el nodo $A$: $Z_A = Z_C \parallel (R_2+Z_L)$[cite: 5]. Exprese $Z_A$ en forma rectangular y polar[cite: 5].

\vspace{0.3cm}
\textbf{b) Análisis fasorial (15 puntos)[cite: 5]} \\
Calcule la tensión $\hat V_A$ y, a partir de ella, determine $\hat V_{out}$ sobre el inductor expresándola en forma rectangular y polar[cite: 5]. Indique si $V_{out}$ adelanta o atrasa a $V_{in}$ y explique qué representa físicamente ese ángulo[cite: 5].

\vspace{0.6cm}
\section*{Problema 2 — Filtro RC con efecto de carga (25 puntos)[cite: 5]}
Una fuente ideal alimenta un filtro pasabajos mediante $R_s = 1\,\text{k}\Omega$[cite: 5]. El filtro contiene $R = 1\,\text{k}\Omega$ en serie y $C = 159.15\,\text{nF}$ entre el nodo de salida y tierra[cite: 5]. Al nodo de salida se conecta una carga resistiva $R_L = 2\,\text{k}\Omega$[cite: 5]. La salida $V_o$ se mide sobre $R_L$, que se encuentra en paralelo con el capacitor[cite: 5].

\begin{center}
    \begin{circuitikz}[scale=1, transform shape, thick]
        \draw (0,0) to[V, v=$\hat V_s$] (0,2.5) to[R, l=$R_s$] (2.5,2.5) to[R, l=$R$] (5,2.5) coordinate(Vo);
        \draw (Vo) to[short, -o] (6.5,2.5) node[right]{$\hat V_o$};
        \draw (Vo) to[C, l_=$C$] (5,0) -- (0,0);
        \draw (6.5,2.5) to[R, l=$R_L$] (6.5,0) -- (5,0);
        \draw (6.5,0) to[short, -o] (6.5,0);
        \draw (5,0) node[ground]{};
    \end{circuitikz}
\end{center}

\textbf{a) Función de transferencia cargada (10 puntos)[cite: 5]} \\
Demuestre analíticamente cómo modifica $R_L$ la función de transferencia[cite: 5]. Parta de $Z_C = 1/(j\omega C)$ y obtenga $H_L(j\omega) = \hat V_o / \hat V_s$[cite: 5]. Simplifique el resultado hasta una expresión de la forma $H_L(j\omega) = K / (1+j\omega\tau)$, determinando $K$ y $\tau$[cite: 5].

\vspace{0.3cm}
\textbf{b) Respuesta frecuencial (15 puntos)[cite: 5]} \\
Evalúe el circuito cargado a $f = 1\,\text{kHz}$ y determine $H_L(j\omega)$ en forma rectangular, $|H_L|$, $\angle H_L$, y la ganancia en dB ($G_{dB} = 20\log_{10}|H_L|$)[cite: 5]. Compare el resultado con el caso $R_L \rightarrow \infty$ y explique físicamente por qué la carga modifica la respuesta[cite: 5].

\end{document}
